\section{Methodology}

\subsection{Overview}

Our goal is to test whether uncertainty probes generalize across LLM families. If uncertainty encoding is architecture-invariant, a probe trained on one model's hidden states should maintain predictive power on another model's states. We operationalize this through a $3 \times 3$ transfer matrix: train three probes (one per model family), evaluate each on all three models' hidden states, and measure the AUROC gap between same-family and cross-family evaluation.

The key insight guiding our design is that if transformers encode uncertainty in a shared geometric structure, linear probes should transfer---and for models with different hidden dimensions, affine alignment should recover the mapping.

\subsection{Semantic Entropy Probe Architecture}

Following \citet{kossen2024semantic}, we train a linear classifier on hidden states:

\textbf{Hidden state extraction.} Given input tokens $x_{1:T}$, we extract the hidden state $h_l \in \mathbb{R}^d$ at layer $l$ and the last token position (Second-Last Token, SLT). We select $l$ at approximately 2/3 depth: layer 21 for 32-layer models (Llama-3, Mistral), layer 18 for the 28-layer model (Qwen-2).

\textbf{Probe training.} We train a logistic regression classifier:
\begin{equation}
P(\text{high-SE} \mid h) = \sigma(w^\top h + b)
\end{equation}
where $w \in \mathbb{R}^d$ and $b \in \mathbb{R}$ are learned parameters. Training labels are binarized semantic entropy: questions with SE above the median are labeled 1 (high uncertainty), others 0. We use sklearn's LogisticRegression with L2 regularization ($C=1.0$) and LBFGS solver.

\textbf{Rationale.} Linear probes are sufficient because prior work shows truthfulness signals are approximately linear in hidden-state space. Logistic regression provides calibrated probabilities and fast training on $\sim$650 examples.

\subsection{Cross-Family Transfer Protocol}

To evaluate whether probes generalize, we construct a transfer matrix:

\begin{enumerate}
\item \textbf{Extract hidden states.} For each of three models $\{m_1, m_2, m_3\}$, extract hidden states from train and validation splits of TruthfulQA (80/20 split, 653/164 questions). Cache states to disk to enable multi-probe evaluation without reloading models.

\item \textbf{Train per-model probes.} Train three SEPs, each on its respective model's training hidden states.

\item \textbf{Evaluate all pairs.} For each (source model, target model) pair:
\begin{itemize}
\item If dimensions match: apply source probe directly to target hidden states
\item If dimensions mismatch: apply affine alignment before evaluation
\end{itemize}

\item \textbf{Compute transfer gaps.} The gap for pair $(i, j)$ is:
\begin{equation}
\text{gap}_{i \to j} = \text{AUROC}_{i \to i} - \text{AUROC}_{i \to j}
\end{equation}
where $\text{AUROC}_{i \to j}$ denotes the AUROC of probe trained on model $i$, evaluated on model $j$ hidden states.
\end{enumerate}

\textbf{Success criterion.} Following our pre-registered threshold, transfer succeeds if all gaps are below 0.10 and mean gap is below 0.05.

\subsection{Affine Alignment for Dimension Mismatch}

Qwen-2-7B has hidden dimension 3584, while Llama-3-8B and Mistral-7B have dimension 4096. Direct probe application fails because the weight vector has wrong size. We address this with affine alignment:

Given paired hidden states from source model (dimension $d_s$) and target model (dimension $d_t$), we learn a mapping:
\begin{equation}
\hat{h}_s = h_t W + b
\end{equation}
where $W \in \mathbb{R}^{d_t \times d_s}$ and $b \in \mathbb{R}^{d_s}$.

\textbf{Fitting.} We solve the least-squares problem on training data:
\begin{equation}
(W^*, b^*) = \arg\min_{W,b} \|H_s - (H_t W + b)\|_F^2
\end{equation}
where $H_s \in \mathbb{R}^{n \times d_s}$ and $H_t \in \mathbb{R}^{n \times d_t}$ are matrices of paired hidden states.

\textbf{Rationale.} Affine alignment is sufficient if the uncertainty subspace is preserved under linear transformation. The least-squares solution is the maximum-likelihood estimator under Gaussian noise. We fit aligners on training split and apply to validation split, preventing overfitting.

\subsection{Implementation Details}

\textbf{Models.} We use instruction-tuned models from HuggingFace: Meta-Llama-3-8B-Instruct, Mistral-7B-Instruct-v0.2, Qwen2-7B-Instruct. All models run in float16 on NVIDIA H100 GPUs.

\textbf{Layer selection.} We extract from layer $\lfloor \text{frac} \times n_{\text{layers}} \rfloor$ with $\text{frac} = 2/3$, following findings that mid-to-upper layers best encode truthfulness. This gives layer 21 for Llama/Mistral (32 layers) and layer 18 for Qwen (28 layers).

\textbf{Caching.} Hidden states are cached as NumPy arrays, reducing GPU memory pressure and enabling rapid iteration on transfer experiments.

\textbf{Reproducibility.} All experiments use seed 42 for train/val splits. Code and cached states will be released upon publication.
